\chapter{Introduction}
\label{ch:intro}

\section{Motivation}
\label{sec:intro_motivation}

Autonomous driving offers the prospect of reshaping road transportation by improving safety, extending access to mobility and enabling a more efficient use of vehicles and infrastructure. The World Health Organization (WHO) estimates that road traffic crashes cause approximately 1.16 million deaths and between 20 and 50 million non-fatal injuries every year worldwide~\cite{who2026}, and reducing the dependence on human driving could help prevent the accidents associated with distraction, fatigue and other human factors. The same reduced dependence could support independent travel for people who are unable to drive, while mobility services based on driverless vehicles could make transportation more accessible and flexible, allowing users to travel according to their needs without owning a vehicle.

Alongside these opportunities, the technology could contribute to a more efficient and sustainable transportation, since smoother driving strategies and an improved coordination between vehicles could reduce energy consumption, emissions and the formation of traffic congestion, while shared fleets could further improve the utilization of the vehicles and reduce the demand for parking, supporting a transition toward mobility provided as a service.

In this perspective, autonomous racing provides a valuable experimental environment for advancing autonomous driving technologies. High speeds, large accelerations, limited reaction times, and close interactions with opponents create demanding conditions in which the performance and robustness of autonomous systems can be evaluated near the limits of vehicle handling. Moreover, by demonstrating the capabilities of autonomous driving to a wider audience, these competitions could help build public confidence in the technology and support its broader acceptance.

Within this setting, this thesis focuses on opponent perception and, in particular, on the estimation of the state of the opponent, critical during close interactions such as overtaking and defending maneuvers. This problem is known in the literature as Target Tracking and consists in reconstructing over time quantities such as position, velocity and acceleration, by associating successive sensor measurements and combining them with the predictions of a vehicle motion model. In the autonomous driving software developed for the Dallara AV-24 and EAV-25 racecars, this process relies on an Extended Kalman Filter (EKF) that integrates the predictions of the motion model with the detections of LiDARs, RaDARs and cameras, and the thesis builds on this framework along two directions. The first is a methodology for the analysis of the measurement pipeline and the characterization of its errors, which lays the foundations for the implementation and testing of an adaptive approach to the estimation of the measurement noise covariance matrices. The second concerns the estimation of the dynamics of the opponent, whose maneuvers entail rapid changes in its motion, which are addressed through the adaptation of the process noise during maneuvers and through non-causal filtering, which uses subsequent observations to refine past state estimates. All the approaches are evaluated on experimental racing data within the existing tracking architecture.

\section{Related Works}
\label{sec:intro_related}

Target tracking has been extensively investigated in the literature. A standard approach to target state estimation is recursive filtering, in which the predictions obtained from a motion model are updated as new sensor observations become available. The Kalman filter provides the classical formulation of this approach for linear systems, as proposed by Kalman~\cite{kalman1960}. For nonlinear problems, widely adopted methods include the Extended Kalman Filter (EKF), the Unscented Kalman Filter (UKF) and the Particle Filter, the last two of which employ more sophisticated methods to propagate and correct the state distribution, as discussed by Wan and van der Merwe~\cite{wan2000}, Julier and Uhlmann~\cite{julier2004} and K\"unsch~\cite{kunsch2013}, but usually at a higher computational cost. The choice among them is therefore typically related to the application. Comparisons of these filters are presented by Tao et al.~\cite{tao2021}, in an automotive context and with different motion models, and by Cui et al.~\cite{cui2005}, in an application to ground target tracking, showing that their characteristics and performance depend on the nonlinearities of the estimation problem, on the noise assumptions and on the representation of the dynamics of the target.

Within this framework, the specification of measurement and process noise covariances is essential because these quantities determine the relative confidence assigned to observations and model predictions. The identification of unknown noise statistics has been extensively studied by Mehra~\cite{mehra1970,mehra1972}, who discusses correlation-based, covariance-matching, maximum-likelihood, and Bayesian approaches. More recently, Zhang et al.~\cite{zhang2020} revisit the identification of noise covariances, examining the conditions under which they can be recovered from data and the constraints required to obtain valid estimates.

When measurement uncertainty varies over time, adaptive filtering provides methods for updating its statistical description from the available data. Fokin and Shchipitsyn~\cite{fokin2009} derive an innovation-based approach that adapts noise covariances using statistics collected over a sliding window. Akhlaghi et al.~\cite{akhlaghi2017} also investigate covariance adjustment using innovations and post-fit residuals. An alternative is offered by variational Bayesian methods, which estimate the state and unknown noise parameters jointly. S\"arkk\"a and Nummenmaa~\cite{sarkka2009} develop this approach for linear state-space models, while S\"arkk\"a and Hartikainen~\cite{sarkka2013} extend it to nonlinear Gaussian filtering and Lan et al.~\cite{lan2023} address the joint estimation of the state and of an unknown process noise covariance. This thesis adopts innovation-based covariance matching for the two LiDAR detection pipelines, recursively updating the measurement noise covariance from the EKF innovations to account for changes in detection quality.

Alongside the uncertainty of the measurements, the accuracy of the model used to predict the state of the opponent is fundamental to obtain a reliable and consistent estimate. Among the possible strategies, the use of more complex and sophisticated models is reviewed by Li and Jilkov~\cite{li2003models} and by Pulford~\cite{pulford2015}, who overview motion models ranging from constant-velocity and constant-acceleration formulations to models describing turning motion, up to more complex formulations such as the one introduced by Singer~\cite{singer1970}, in which the acceleration of the target is represented as a correlated stochastic process.

Another widely adopted line of research to account for the uncertainty of the model is represented by multiple-model methods, which consider several motion hypotheses simultaneously. A bank of filters, each based on a different dynamic model, runs in parallel, and the estimate is obtained by combining their outputs according to the probability that each model describes the current motion of the target. The Interacting Multiple Model (IMM) algorithm introduced by Blom and Bar-Shalom~\cite{blom1988} is the most established example, in which the filters interact at each step through the mixing of their estimates, according to a Markov model of the switching between the modes. Li and Jilkov~\cite{li2005multiple} provide a comprehensive review of this family of methods, Mazor et al.~\cite{mazor1998} survey the application of the IMM to target tracking, while Kolat et al.~\cite{kolat2022} investigate its application to vehicle maneuver classification.

A complementary family of approaches uses explicit maneuver detection to modify the behavior of the estimator. Li and Jilkov~\cite{li2002decision} and Pulford~\cite{pulford2015} review decision-based methods, including equivalent-noise techniques, input detection and estimation, and model switching. Examples of this approach are the variable dimension filter of Bar-Shalom and Birmiwal~\cite{barshalom1982}, which switches to a higher-order motion model when a maneuver is detected, and the adaptive estimator of Moose~\cite{moose1975}, which models the maneuver as a semi-Markov process within a Bayesian estimation scheme. Related methods include strong tracking filters, reviewed by Crouse~\cite{crouse2019}, which modify the predicted uncertainty to improve robustness to model mismatch. In this thesis, a decision-based strategy is implemented through a jerk-based maneuver detection logic that drives a feed-forward adaptation of the process noise covariance. An auxiliary EKF provides the acceleration estimate from which jerk is computed, while a hysteresis-based logic switches the main filter between nominal and inflated process noise levels to improve responsiveness during abrupt maneuvers.

Moving beyond the hypothesis of causal filtering, i.e.\ the prediction of the motion model followed by the correction with the current measurement, a part of the literature is dedicated to non-causal approaches, which improve the estimates of past states by incorporating the observations acquired after the instant being estimated. For linear systems, the Rauch--Tung--Striebel (RTS) smoother performs a backward recursion following a forward filtering pass, as shown by Rauch et al.~\cite{rauch1965}, while Fraser and Potter~\cite{fraser1969} formulate smoothing through the combination of forward and backward filters. Wall et al.~\cite{wall1981} further investigate the fixed-interval smoothing problem. Fixed-lag formulations, studied by Moore~\cite{moore1973}, limit the subsequent information to a prescribed window, making the trade-off between estimation accuracy and delay explicit. S\"arkk\"a~\cite{sarkka2013book} overviews the application of the RTS smoother in different contexts, from linear systems to nonlinear ones, through its counterparts based on the EKF and on the UKF, up to particle smoothers. In our work, the RTS formulation is implemented over a sliding window of fixed duration. As new measurements are processed, the window advances and the backward recursion refines the past state estimates within it, incorporating the newly available information and enabling continuous processing of the measurement stream with a prescribed lag.

\section{Contributions and Results}
\label{sec:intro_contributions}

The first part of this thesis addresses the characterization of the measurement noise, with the aim of obtaining covariance values that reflect the actual uncertainty of the detections, so that each measurement is weighted consistently with its accuracy in the correction step of the EKF. The analysis considers the four detection sources of the tracker, i.e.\ LiDAR Clustering, LiDAR PointPillars, RaDAR and camera, and is based on experimental racing data, in which the detections are compared with the reconstructed ground truth of the opponent.

The errors are first cleaned of the deterministic contributions related to the reference conventions of the ground truth and to the extrinsic calibration of the sensors, and then of the systematic and repeatable components related to the operating principle of each sensor and detection pipeline. The residual errors are subsequently analyzed in the time and frequency domains through their normalized autocorrelation and power spectral density, and the latter analysis shows the presence of a slowly varying, time-correlated component, on which a white noise is superimposed. This component is modeled as a first-order Gauss--Markov process, characterized by its time constant and variance, which is used offline to reconstruct and remove it from the residuals, so that the white component on which the covariance is tuned can be isolated.

The distribution of the decorrelated residuals remains leptokurtic, particularly for LiDAR Clustering, reflecting the strong dependence of measurement quality on the observation conditions. Most samples form a dense central core, corresponding to small errors when the sensor observes the opponent clearly, while the heavy tails represent more uncertain measurements obtained under less favorable conditions. This variability motivates an innovation-based covariance-matching approach for the two LiDAR pipelines, which adapts the measurement noise covariance recursively to changes in detection quality. Averaging the resulting time-varying Gaussian models produces a mixture capable of representing both the narrow core and the heavier tails.

The consistency of this approximation is assessed against the empirical distribution of the decorrelated residuals, using the nominal static tuning and a Gaussian fitted to the white component as benchmarks. Relative to the latter, the distribution coverage for LiDAR Clustering increases from \SI{76}{\percent} to \SI{92}{\percent} longitudinally and from \SI{75}{\percent} to \SI{88}{\percent} laterally. For PointPillars, coverage increases from \SI{85}{\percent} to \SI{90}{\percent} and from \SI{84}{\percent} to \SI{90}{\percent}, respectively.

Beyond the quantitative results of the covariance matching, a further outcome of this first part is the definition of a systematic methodology for the analysis of the errors and trends of the sensors and for the selection of the measurement noise covariance. This makes it possible, whenever the detection pipelines are modified in the future, to repeat the same procedure in order to identify and quantify the changes in their error characteristics and to update the noise model and its tuning consistently, without redeveloping the analysis.

The second part examines the process noise covariance and its influence on the estimation of the opponent's dynamics, particularly during aggressive braking and acceleration maneuvers. A feed-forward adaptation is developed to increase the responsiveness of the tracker when a maneuver occurs. An auxiliary EKF provides an acceleration estimate from which the jerk is derived, while a hysteresis-based detection logic switches the process noise of the main filter between nominal and inflated values. Experimental results show that this strategy improves the response during detected maneuvers while retaining smoother estimates during steady motion, although the quality of the jerk estimate influences the timing and reliability of the activation.

The limitations of this approach motivate the investigation of non-causal filtering techniques to improve the reconstruction of the opponent's dynamics. In particular, a Rauch--Tung--Striebel (RTS) smoother is implemented with a fixed-lag sliding window, using subsequent observations to refine past state estimates. The influence of the window length is investigated in terms of both estimation accuracy and temporal alignment of the reconstructed maneuvers with the ground truth. With a window of \SI{300}{ms}, the speed RMSE is reduced by approximately \SI{31}{\percent} when all sensors are available and by \SI{59}{\percent} when RaDAR measurements are excluded. The latter result demonstrates the benefit of non-causal estimation when speed must be inferred from position measurements alone.

Finally, the RTS smoother is combined with the feed-forward process noise adaptation to address the limitations of the maneuver detection strategy. Although the feed-forward adaptation can react only after the jerk-based detector activates, the smoother uses subsequent information to refine the reconstructed onset retrospectively and attenuate fluctuations caused by delayed or spurious activations. The combination yields the lowest estimation errors among the configurations considered: relative to the causal feed-forward estimate, the speed RMSE decreases from \SI{1.069}{m/s} to \SI{0.356}{m/s}, a reduction of \SI{66.7}{\percent}, while the acceleration RMSE decreases from \SI{3.589}{m/s^2} to \SI{1.786}{m/s^2}, a reduction of \SI{50.2}{\percent}.

\section{Thesis Structure}
\label{sec:intro_structure}

The remainder of the thesis is organized as follows.

\begin{itemize}
\item Chapter~\ref{ch:setup} introduces the autonomous racing competitions, the vehicles employed in this work, their sensing equipment and the architecture of the autonomous driving software.

\item Chapter~\ref{ch:tracking} describes the target tracking pipeline, including the Extended Kalman Filter, the motion and measurement models, data association and track management.

\item Chapter~\ref{ch:adaptive_r} presents the experimental characterization of the measurement errors, their statistical and temporal properties, and the development and validation of an adaptive measurement noise covariance.

\item Chapter~\ref{ch:qff} investigates the effect of the process noise on the estimation of the opponent's dynamics and introduces the maneuver detection and feed-forward adaptation strategy.

\item Chapter~\ref{ch:rts} develops a non-causal state estimator based on a fixed-lag RTS smoother, presents its formulation and implementation, performs a sensitivity analysis of its tuning parameters and evaluates its combination with the feed-forward process noise.

\item Chapter~\ref{ch:conclusions} summarizes the main findings and outlines possible future developments.
\end{itemize}
